\documentclass[11pt]{article}
\usepackage[left=0.3in,right=0.6in,top=1in,bottom=1in]{geometry}

%others
\newtheorem{thm}{Theorem}[section]
\newtheorem{lemma}[thm]{Lemma}
\newtheorem{sublemma}[thm]{Sub-Lemma}
\newtheorem{prop}[thm]{Proposition}
\newtheorem{corollary}[thm]{Corollary}

%% Theorem style - definition
\newtheorem{defn}[thm]{Definition}
\newtheorem{example}[thm]{Example}
\newtheorem{remark}[thm]{Remark}
\newtheorem{conj}[thm]{Conjecture}
\newtheorem{examples}[thm]{Examples}
\newtheorem{obs}[thm]{Observation}
\newtheorem{claim}[thm]{Claim}
\newtheorem{assumption}[thm]{Assumption}
\newtheorem{notation}[thm]{Notation}
\newtheorem{question}[thm]{Question}
\newtheorem{construction}[thm]{Construction}
\newtheorem{outline}[thm]{Outline}
\newtheorem{assumptions}[thm]{Standing Assumptions}

\usepackage{amsmath, amssymb, amsfonts}
%\usepackage{mathabx}
\usepackage{graphicx, overpic}
%\usepackage{comment}
\usepackage{color}
\usepackage{enumerate}
\usepackage{multicol}
%\usepackage{xypic}
\usepackage{mathrsfs}
\usepackage{mathabx}
\usepackage{proof}
\usepackage[T1]{fontenc}


% The [listings] package is very helpful for typesetting code.  It is aware of
% many languages---i.e., it knows what the keywords are and things like that,
% and can adjust how they are displayed.  It has a lot of features for
% controlling the display, but we'll just use a few here.  [listtings] lets
% you define "styles", which switch between different settings for different
% languages.
\def\utt#1{\underline{\smash{\texttt{#1}}}}
\usepackage{listings}
  % This is a simple style for displaying OCaml code.
  \lstdefinestyle{ocamlcode}{
    language=[Objective]Caml,
	  basicstyle=\normalfont,
	  keywordstyle=\utt,
	  morekeywords={int, bool, real, char, string, true, false, not,
	  	div, mod, list},
	  identifierstyle=\texttt,
	  commentstyle=\texttt,
	  stringstyle=\texttt,
	  tabsize=2,
	  showstringspaces=false,
	  mathescape=true,
	  alsoletter={0,1,2,3,4,5,6,7,8,9,-,<,>,'},
  }
  % This is a simple style for displaying C code.
  \lstdefinestyle{ccode}{
    language=C,
    keywordstyle=\normalfont\textttunder,
    keywordstyle=[2]\normalfont\textt,
    morekeywords={int,bool},
    identifierstyle=\normalfont\texttt,
    commentstyle=\normalfont\texttt,
    stringstyle=\normalfont\texttt,
    tabsize=4,
    showstringspaces=false,
    mathescape=true,
    alsoletter={0,1,2,3,4,5,6,7,8,9,0,<,>,=,+,-,*,/,|,!,),(,&,\%,[,]},
  }
  % This macro enables a short form for typesetting code.
  \lstMakeShortInline|

%%% Used for typesetting inference rules and tree-style derivations; provides
%%% the \inferrule macro.
\usepackage{mathpartir}

%%% Used for rotated figures
\usepackage{rotating}

%%% Used for typesetting derivations.
\usepackage{fitch}
    \newcommand{\der}[3][{}]{\hypo{#1}{#2} #3 }
    \newcommand{\subder}[3][{}]{\open\hypo{#1}{#2} #3 \close}
    \newcommand{\tbc}[1][{}]{\open\have{#1}{\cdots}\close}
    \newcommand{\axiom}[2][{}]{\subder[#1]{#2}{}}

%%% Our languages
\newcommand{\Ocm}{\mbox{\textrm{OCaml}-}}
\newcommand{\Cm}{\mbox{\textrm{C}-}}

% Writing inference rules often involves mixing mathematics with listings.
% For example, a generic expression should be e, typeset as a mathematical
% symbol.  So a generic Ocaml- let expression should be typeset as
% \lstinline!let $x$ = $e'$ in $e$!.  Unfortunately, the spacing can get very
% messed up if you do this in arguments to macros, and you need to put some
% spaces back in.  So I often write macros like these that use \, to insert
% space and use these macros when needed.

\newcommand{\Ocmlet}[3]{\lstinline[style=ocamlcode]!let $\,#1\,$ = $\,#2\,$ in $\,#3$!}

%%% Environments.
\newcommand{\set}[1]{\left\{#1\right\}}
\newcommand{\env}[1]{\set{#1}}
\newcommand{\bind}[2]{{#1}\mapsto{#2}}
\newcommand{\extend}[3]{#1\env{\bind{#2}{#3}}}
\newcommand{\emptyenv}{\_}
\newcommand{\wild}{\_}
\newcommand{\rhos}{\rho s}

%%% Judgments

% \expjudge[p] ρs e v = ρs ├_p e ↓ v
\newcommand{\expjudge}[4][{}]{{#2}\vdash_{#1}{#3}\Downarrow{#4}}
% \stmjudge[p] s ρs ρs' = s ├_p ρs → ρs'
\newcommand{\stmjudge}[4][{}]{{#2}\vdash_{#1}{#3}\Rightarrow{#4}}


%%% Title, author, date, 
\title{CS324 - Project 2}
\author{Jeramiah Trout, Ruishi Wang, Mark Zheng}
\date{\today}


\begin{document}

\maketitle
    
\noindent Task 2 - Operational Semantics

\textbf{Example of a program that is not safe}:
\begin{center}
    \begin{verbatim}
      bool b;
      int x;
      int y;
      fscanf(stdin_hi, "%b", &b);
      if (b) {
          fscanf(stdin_lo, "%d", &x);
      }
      fscanf(stdin_lo, "%d", &y);
      fprintf(stdout_lo, "%d", y);
    \end{verbatim}
\end{center}

\noindent By Definition 2 in handout, a program is safe when public inputs are the same, public outputs are also same even 
secrete inputs can be different. 

To show our program above is not saft, then, we need to run the program twice, with two different high security input. We can 
do this in our program by changing boolean value b from true to false, then compare the output result. 

Looking at the derivation on the next page, we have the following result

\[
  chs_0 = \langle [1,2], [\,], [true], [\,]\rangle
  \qquad
  chs_1 = \langle [1,2], [\,], [false], [\,]\rangle
  \]
  \[
  chs_0' = \langle [\,], [2^L], [\,], [\,]\rangle
  \qquad
  chs_1' = \langle [2], [1^L], [\,], [\,]\rangle
  \]
  \[
  chs_0 \sim chs_1
  \qquad\text{but}\qquad
  chs_0' \not\sim chs_1'
  \]

\begin{figure}
  \centering
  \resizebox{\linewidth}{!}{$
  \begin{nd}
  \der{
    ss \vdash_p \_,\, chs_0 \Rightarrow \eta_0,\, chs_0'
    \qquad
    ss \vdash_p \_,\, chs_1 \Rightarrow \eta_1,\, chs_1'
  }
    % ---------- First run: b = true ----------
    \subder{
      ss \vdash_p \_,\, chs_0 \Rightarrow \eta_0,\, chs_0'
    }{
      \axiom{
        \texttt{bool b; int x; int y;}
        \vdash_p \_
        \Rightarrow
        \{b\mapsto ?^L,\;x\mapsto ?^L,\;y\mapsto ?^L\}
      }
  
      \axiom{
        \texttt{fscanf}(stdin_{hi},\dots,\&b)
        \vdash_p
        \{b\mapsto ?^L,\;x\mapsto ?^L,\;y\mapsto ?^L\},\, chs_0
        \Rightarrow
        \{b\mapsto true^H,\;x\mapsto ?^L,\;y\mapsto ?^L\},\, chs_0^1
      }
  
      \subder{
        \texttt{if}(b)\{
        \texttt{fscanf}(stdin_{lo},\dots,\&x);
        \}
        \vdash_p
        \{b\mapsto true^H,\;x\mapsto ?^L,\;y\mapsto ?^L\},\, chs_0^1
        \Rightarrow
        \{b\mapsto true^H,\;x\mapsto 1^H,\;y\mapsto ?^L\},\, chs_0^2
      }{
        \axiom{
          \{b\mapsto true^H,\;x\mapsto ?^L,\;y\mapsto ?^L\}
          \vdash_p b \Downarrow true^H
        }
  
        \axiom{
          \texttt{fscanf}(stdin_{lo},\dots,\&x)
          \vdash_{p,H}
          \{b\mapsto true^H,\;x\mapsto ?^L,\;y\mapsto ?^L\},\, chs_0^1
          \Rightarrow
          \{b\mapsto true^H,\;x\mapsto 1^{L\sqcup H},\;y\mapsto ?^L\},\, chs_0^2
        }
      }
  
      \axiom{
        \texttt{fscanf}(stdin_{lo},\dots,\&y)
        \vdash_p
        \{b\mapsto true^H,\;x\mapsto 1^H,\;y\mapsto ?^L\},\, chs_0^2
        \Rightarrow
        \{b\mapsto true^H,\;x\mapsto 1^H,\;y\mapsto 2^L\},\, chs_0^3
      }
  
      \axiom{
        \texttt{fprintf}(stdout_{lo},\dots,y)
        \vdash_p
        \{b\mapsto true^H,\;x\mapsto 1^H,\;y\mapsto 2^L\},\, chs_0^3
        \Rightarrow
        None,\, chs_0'
      }
    }
  
    % ---------- Second run: b = false ----------
    \subder{
      ss \vdash_p \_,\, chs_1 \Rightarrow \eta_1,\, chs_1'
    }{
      \axiom{
        \texttt{bool b; int x; int y;}
        \vdash_p \_
        \Rightarrow
        \{b\mapsto ?^L,\;x\mapsto ?^L,\;y\mapsto ?^L\}
      }
  
      \axiom{
        \texttt{fscanf}(stdin_{hi},\dots,\&b)
        \vdash_p
        \{b\mapsto ?^L,\;x\mapsto ?^L,\;y\mapsto ?^L\},\, chs_1
        \Rightarrow
        \{b\mapsto false^H,\;x\mapsto ?^L,\;y\mapsto ?^L\},\, chs_1^1
      }
  
      \subder{
        \texttt{if}(b)\{
        \texttt{fscanf}(stdin_{lo},\dots,\&x);
        \}
        \vdash_p
        \{b\mapsto false^H,\;x\mapsto ?^L,\;y\mapsto ?^L\},\, chs_1^1
        \Rightarrow
        \{b\mapsto false^H,\;x\mapsto ?^L,\;y\mapsto ?^L\},\, chs_1^1
      }{
        \axiom{
          \{b\mapsto false^H,\;x\mapsto ?^L,\;y\mapsto ?^L\}
          \vdash_p b \Downarrow false^H
        }
      }
  
      \axiom{
        \texttt{fscanf}(stdin_{lo},\dots,\&y)
        \vdash_p
        \{b\mapsto false^H,\;x\mapsto ?^L,\;y\mapsto ?^L\},\, chs_1^1
        \Rightarrow
        \{b\mapsto false^H,\;x\mapsto ?^L,\;y\mapsto 1^L\},\, chs_1^2
      }
  
      \axiom{
        \texttt{fprintf}(stdout_{lo},\dots,y)
        \vdash_p
        \{b\mapsto false^H,\;x\mapsto ?^L,\;y\mapsto 1^L\},\, chs_1^2
        \Rightarrow
        None,\, chs_1'
      }
    }
  \end{nd}
  $}
  \end{figure}
\end{document}
